\cotefile{157-5}
\scanpage{437}
\mspage{349}
% slanted
\margin{À rejeter plus bas, après 4.13.31, qui \unsure{est|donne} la première \unsure{application} importante de 4.13.22, alors que 4.13.28 se prouve sans avoir à utiliser 4.13.22 \unsure{ou|ni} ses variantes.}
\uline{Proposition} 4.13.28. \unsure{Soient|\deleted{Soient}} $M,\pi,\Phi$ comme dans \deleted{\unsure{4}} 4.13.23. Supposons de plus que \deleted{\unsure{$\Phi$}} $J\in\Phi$ implique que $J$ est \uline{filtrante}. Alors l'ens $\widetilde{\Sigma}_{\pi}^{\Phi}$ \unsure{des|\deleted{des}} sous-ensembles $\Phi$-$\pi$-accessibles de $M$ est stable par \deleted{\unsure{réunions}} réunions finies. \deleted{Supposons} \inserted{Plus \unsure{gén$^{\rm t}$},} soit $\lambda$ un cardinal \inserted{$<\pi$} tel que \deleted{\unsure{$\Phi$}} $\forall\,J\in\Phi$, $J$ soit grde devant $\lambda$. Alors l'ens $\widetilde{\Sigma}_{\pi}^{\Phi}$ est stable par réunions \deleted{\unsure{index}} $(S_{\unsure{r|i}})_{r\in A}$ \unsure{indexées|indiciées} par des ens. $A$ de cardinal $\leq\lambda$.

\uline{Corollaire} 4.13.29. L'ens des parties accessibles de $\mathrm{Ob}\,M$ est stable par \inserted{``petites''} réunions --- i.e. si $(S_r)_{r\in A}$ est une \inserted{\unsure{petite}} famille de parties accessibles de $\mathrm{Ob}\,M$, alors $\bigcup_{r\in A}S_r$ est aussi accessible.

La démonstration de 4.13.28 repose sur le lemme suivant (je n'explicite pas les démonstrations):
% archive sheet number: 436
